\begin{Answer}{2.4}
$2d+1$;
$c+d+1$;
even
\end{Answer}
\begin{Answer}{2.6}
$2n$;
$2o+1$;
some integers $n$ and $o$;
$4no+2n=2(2no+n)$ or $2(n(2o+1))$.  (Your steps might vary slightly, but you should end up with
either $2(2no+n)$ or  $2(n(2o+1))$ in the final step);
$2no + 1$  or $n(2o+1)$;
`an even integer' or `even'.

\end{Answer}
\begin{Answer}{2.7}
Let $a$ and $b$ be even integers.  Then $a=2m$ and $b=2n$ for some integers $m$ and $n$.
Their product is $ab=(2m)(2n)=2(2mn)$ which is even since $2mn$ is an integer. {\em Notice that we used two different letters here!
You cannot assume $a=2n$ and $b=2n$  because then you are
assuming that $a=b$ whether or not you realize it!}
\end{Answer}
\begin{Answer}{2.8}
Here are my comments on the proof.
\begin{itemize}
\item
The first sentence is phrased weird--we are not letting $a$ and $b$ be odd {\em by} the definition of odd.
We are {\em using} the definition.
\item
It does not state that $n$ and $q$ need to be integers.
\item
Although it is not incorrect, using $n$ and $q$ is just weird in this context.
It is customary to use adjacent letters, like $n$ and $m$, or $q$ and $r$.
\item
Given the above problems, I would rephrase the first sentence as
{\em `Let $a$ and $b$ be an odd numbers.
Then $a=2n+1$ and $b=2m+1$ for some integers $n$ and $m$.'}
\item
There is an algebra mistake.  The product should be $2(2nq+q+n)$.
\item
If you replace $2nq+1$ with $2nq+q+n$ (twice) in the last sentence (see the previous item)
it would be a perfect finish to the proof.
\end{itemize}
\end{Answer}
\begin{Answer}{2.9}
Hopefully it is clear to you that the proof {\em can't} be correct since the sum of an even and an odd number is odd, not even.
The algebra is correct.
{\em The problem is that $n+m+1/2$ is not an integer.}
In order to be even, a number must be expressed
in the form $2k$ {\bf\em where $k$ is an integer}.  {\em Any} number can be written as $2x$ if we don't require that $x$ be an integer,
so you {\em cannot} say that a number is even because it is of the form $2x$ unless $x$ is an integer.
\end{Answer}
\begin{Answer}{2.13}
$a$ an integer; $(3x + 2)$; $(5x - 7)$;$7$; $7$ divides $15x^2 - 11x - 14$.
\end{Answer}
\begin{Answer}{2.15}
This proof is correct.  Not all of the Evaluate problems have an error!
\end{Answer}
\begin{Answer}{2.17}
The number $2$ is positive and even but is clearly not composite since it is prime.
Since the statement is false the proof must be incorrect.
So where is the error?  It is in the final statement.  Although $a$ can be written
as the product of $2$ and $k$, what if $k=1$ (that is, $a=2$).  In that case
we have not demonstrated that $a$ has a factor other than $a$ or $1$, so we can't be sure
that it is composite.
\end{Answer}
\begin{Answer}{2.18}
If you didn't get, try this hint before reading the rest of the solution:
 Assume $a$ is an even number other than $2$ and prove that $a$ is composite.\\
Let $a>2$ be an even integer.  Then $a=2k$ for some integer $k$.  Since $a\neq 2$,
$a$ has a factor other than $a$ or $1$.  Therefore $a$ is not prime.
Therefore $2$ is the only even prime number.
\end{Answer}
\begin{Answer}{2.19}
It was O.K. because according to the definition of prime, only positive integers can be prime.
Therefore we only needed to consider positive even integers.
\end{Answer}
\begin{Answer}{2.23}
This one has a combination of two subtle errors.
First of all, if $a|c$ and $b|c$, that does not necessarily imply that $ab|c$.
For instance, $6|12$ and $4|12$, but it should be clear that $6\cdot 4\nmid 12$.
Second, what if $a=b$?
We'll see how to fix the proof in the next example.
\end{Answer}
\begin{Answer}{2.25}
Since $n$ is not a perfect square, we know that $a\neq b$.
Therefore  $a<b$ or $b<a$.  Since $a$ and $b$ are just labels for
two factors of $n$, it doesn't matter which one is larger.  So we can
just assume $a$ is the smaller one without any loss of generality.
By definition of composite, we know that $a>1$.
Finally, it should be pretty clear that $b<n-1$ since if $b=n-1$, then
$n=ab=a(n-1)\geq 2(n-1)=2n-2=n+(n-2)>n$ since $n>4$. But clearly $n>n$ is impossible.
\end{Answer}
\begin{Answer}{2.26}
We assumed that $n=a^2>4$, so clearly $a>2$.
\end{Answer}
\begin{Answer}{2.28}
\begin{enumerate}
\item {\bf Experiment}.  If you aren't sure what to do, don't be afraid to try things.
\item {\bf Read Examples}. But don't just read.  Make sure you {\em understand} them.
\item {\bf Practice}.  It makes perfect!
\end{enumerate}
\end{Answer}
\begin{Answer}{2.32}
Only when you read \xkcd{} and you don't laugh.
\end{Answer}
\begin{Answer}{2.33}
If you build it and they don't come, the proposition is false.  This is the only case where it is false.
To see this, notice that if you build it and they do come, it is true.  If you don't build it, then it
doesn't matter whether or not they come--it is true.
\end{Answer}
\begin{Answer}{2.35}
If you don't know a programming language, then you don't know Java.
\end{Answer}
\begin{Answer}{2.37}
true; $\neg p$; false; $p$; $p$ is true; $q$ is false (the last two can be in either order).
\end{Answer}
\begin{Answer}{2.39}
If you don't know Java, then you don't know a programming language.
\end{Answer}
\begin{Answer}{2.40}
They are {\em not} equivalent.  Since Java is a programming language, the proposition seems obviously true.
However, what if someone knows C++ but not Java?  Then they know a programming language but they don't
know Java. Thus, the inverse is false.  Since one is true and the other is false,
the proposition and its inverse are clearly not equivalent.
\end{Answer}
\begin{Answer}{2.42}
If you know a programming language, then you know Java.
\end{Answer}
\begin{Answer}{2.43}
They are {\em not} equivalent.  Since Java is a programming language, the proposition seems obviously true.
However, what if someone knows C++ but not Java?  Then they know a programming language but they don't
know Java. Thus, the converse is false.  Since one is true and the other is false,
the proposition and its converse are clearly not equivalent.
\end{Answer}
\begin{Answer}{2.46}
(a)
The implication states that if I get to watch ``The Army of Darkness'' that I will be happy.  However,
it doesn't say that it is the only thing that will make me happy.
For instance, if I get to see ``Iron Man'' instead, that would also make me happy.
Thus, the inverse statement is false.\\
(b)
I will use fact that $p\rightarrow q$ is true unless $p$ is true and $q$ is false.
The implication is true unless I watch ``The Army of Darkness'' and I am not happy.
The contrapositive is ``If I am not happy, then I didn't get to watch `The Army of Darkness.' ''
This is true unless I am not happy and I watched ``The Army of Darkness.''
Since this is exactly the same cases in which the implication are true, the implication
and its contrapositive are equivalent.
\end{Answer}
\begin{Answer}{2.49}
$ \sqrt{35}$; $10\sqrt{35}$; $3481 \geq 3500$; {\em nonsense} or {\em false} or {\em a contradiction}.
\end{Answer}
\begin{Answer}{2.50}
\begin{pfevalenum}
\item
Here are my comments on this proof:
\begin{itemize}
\item
It is proving the wrong thing.  This proves that the product of an even number and
an odd number is even.  But it doesn't even do that quite correctly as we will see next.
\item
The first sentence is phrased weird--we are not letting $a$ be even {\em by} the definition of even.
We are {\em using} the definition.
\item
It does not state that $n$ and $q$ need to be integers.
\item
Although it is not incorrect, using $n$ and $q$ is just weird.
It is customary to use adjacent letters, like $n$ and $m$, or $q$ and $r$.
\item
Given the above problems, I would rephrase the first sentence as
{\em `Let $a$ be an even number and $b$ be an odd number.
Then $a=2n$ and $b=2m+1$ for some integers $n$ and $m$.'}
\item
There is an algebra mistake.  The product should be $2(2nq+n)$.
\item
The last sentence is actually perfect (again, except for the fact that it isn't proving the right thing).
\end{itemize}
\item
This proof is incorrect.  It actually proves the {\em converse} of the statement.  (We'll learn more about
converse statements later.)
In other words, it proves that if at least of one of $a$ or $b$ is even, then $ab$ is even.
This is {\em not} the same thing.
It is a pretty good proof
of the wrong thing, but it can be improved in at least 4 ways.
\begin{itemize}
\item
It defines $a$ and $b$ but never really uses them.
They should be used at the beginning of the algebra steps (i.e. $a\cdot b=\cdots$) to make
it clear that the algebra is related to the product of these two numbers.
\item
It needs to state that $k$ and $x$ are integers.
\item
As above, using $k$ and $x$ is weird (but not wrong).  It would be better to use $k$ and $l$, or $x$ and $y$.
\item
It needs a few words to bring the steps together.
In particular, sentences should not generally begin with algebra.
\end{itemize}
Taking into account these things, the second part could be rephrased as follows.
\begin{quote}
Let $a=2n$ and $b=2m+1$, where $n$ and $m$ are integers.
Then $ab=(2n)(2m+1)=4nm+2n=2(2nm+n)$, which is even since $2nm+n$ is an integer.
\end{quote}
\item
This proof is correct.
\end{pfevalenum}
\end{Answer}
\begin{Answer}{2.54}
$(1,2,3)$, $(1,3,2)$, $(2,1,3)$, $(2,3,1)$, $(3,1,2)$, and $(3,2,1)$.
\end{Answer}
\begin{Answer}{2.56}
Since it wasn't obvious how to do a direct proof of the fact, proof by contradiction seemed like
the way to go.  So we begin by assuming what we want to prove (that the product is even) is false.
The short answer: {\em Because contradiction proofs generally begin by assuming the negation of what
you want to prove.}
\end{Answer}
\begin{Answer}{2.57}
The proof gives the justification for this, but you may have to think about it
for it to entirely sink in.
Consider carefully the definition of $S$:
$S = (a_1 - 1) + (a_2 - 2) + \cdots + (a_n - n).$
Notice it adds and subtracts terms.  If $S=0$, then the amount added and subtracted
must be the same.  And if you think about it for a few minutes, especially in light
of the justification given in the proof, you should see why.
If you can't see it right away, go back to how the $a_k$'s are defined and think a little more.
If you get totally stuck, try an example with $n=3$ or $4$.
\end{Answer}
\begin{Answer}{2.60}
Because $a^2=a\cdot a$, so to list the factors of $a^2$ you can list the factors of $a$ twice.
Thus, $a^2$ has twice as many factors as $a$, so it must be an even number.
\end{Answer}
\begin{Answer}{2.63}
(1) No. (2) Yes. (3) No. (4) No.
(5) Statements of the form ``$p$ implies $q$'' are false
precisely when $p$ is true and $q$ is false.
(6)  No.  Whether or not you are 21, you aren't breaking the rule.
(7) No.  If $p$ is false, whether or not $q$ is true or false doesn't matter--the statement is true.
Let's consider the previous question--if you do not drink alcohol, you are following the rule
regardless of whether or not the statement ``you are 21'' is true or false.
\end{Answer}
\begin{Answer}{2.64}
$a > b$; $\frac{a - b}{2}$; $\frac{a - b}{2}$; $b+\frac{a}{2}-\frac{b}{2}=\frac{a}{2}+\frac{b}{2}$;
subtract $\frac{a}{2}$ from both sides and multiple both sides by $2$;
$a > b$; contradiction; $a\leq b$.
\end{Answer}
\begin{Answer}{2.66}
$a\left(\frac{p}{q}\right)^2 + b\left(\frac{p}{q}\right) + c$;
multiple both sides by $q^2$;
odd;
0;
$ap^2 + bpq$ is even and $cq^2$ is odd, so  $ap^2 + bpq + cq^2$ is odd;
$bpq + cq^2$ is even and $ap^2$ is odd, so  $ap^2 + bpq + cq^2$ is odd;
$ax^2 + bx + c = 0$ does not have a rational solution if $a$, $b$, and $c$ are odd.
\end{Answer}
\begin{Answer}{2.70}
\begin{pfevalenum}
\item
This is attempting to prove the {\em converse}, not the contrapositive.
Since the converse of a statement is not equivalent to the original statement, this is not a valid proof.
Further, the proof contains an algebra mistake.
Finally, it uses the property that the sum of two even integers is even.  Although this is true, the problem
specifically asked to prove it using the definition of even/odd.
\item
This proof starts out correctly by using the contrapositive statement and the definition of odd.
Unfortunately, the writer claims that $\displaystyle 5(\frac{6}{5}k+1)$ is `clearly odd.'
This is not at all clear.  What about this number makes it odd?  Is it expressed as $2a+1$ for
some integer $a$?  No.
Even worse, there is a fraction in it, obscuring the fact that the number is even an integer.
\item
This proof is {\em really close}.  The only problem is that we don't know that $6k+5$ is odd
{\em using the definition of odd}.  All the writer needed to do is take their algebra a little
further to obtain $2(3k+2)+1$, which is odd by the definition of odd since $3k+2$ is an integer.
\end{pfevalenum}
\end{Answer}
\begin{Answer}{2.76}
Answers will vary greatly, but one proof is: $3$ and $5$ are prime but $3+5=8=2^3$ is clearly not prime.
\end{Answer}
\begin{Answer}{2.79}
  $2s$ is a power of two that is in the closed interval.;
$2^r=2\cdot 2^{r-1}< 2s<2\cdot 2^r=2^{r+1}$, so
$s< 2^r < 2s < 2^{r+1}$, and so the interval $[s,2s]$ contains $2^r$,
a power of $2$.
\end{Answer}
\begin{Answer}{2.80}
Because these statements are contrapositives of each other.  In other words, they are equivalent.
Therefore you can prove either form of the statement.
\end{Answer}
\begin{Answer}{2.82}
If $x$ is odd, then $x=2k+1$ for some integer $k$.
Then $x+20=2k+1+20=2(k+10)+1$, which is odd since $k+10$ is an integer.
If $x+20$ is odd, then $x+20=2k+1$ for some integer $k$.  Then $x=(x+20)-20=2k+1-20=2(k-10)+1$,
which is odd since $k-10$ is an integer.
Therefore $x$ is odd iff $x+20$ is odd.
\end{Answer}
\begin{Answer}{2.83}
If $x$ is odd, then $x=2k+1$ for some integer $k$.
Then $x+20=2k+1+20=2(k+10)+1$, which is odd since $k+10$ is an integer.
If $x$ is even, then $x=2k$ for some integer $k$.
Then $x+20=2k+20=2(k+10)$.  Since $k+10$ is an integer, then
$x+20$ is even. Therefore $x$ is odd iff $x+20$ is odd.
\end{Answer}
\begin{Answer}{2.84}
$p$ implies $q$;
$q$ implies $p$;
$p$ implies $q$;
$\neg p$ implies $\neg q$
\end{Answer}
\begin{Answer}{2.85}
\begin{pfevalenum}
\item
For the forward direction, they didn't use the definition of odd.  Otherwise, that part is fine.
For the backward direction, their proof is nonsense.  They {\em assumed} that $x=2k+1$ when they wrote
$(2k+1)-4$ in the second sentence.  This need to be {\em proven}.
\item
For the forward direction, they didn't specify that $k$ was an integer.  Otherwise it is correct.
The second part of the proof is {\em not} proving the converse.  It is proving the forward direction
a second time using a proof by contraposition.  In other words, this proof just proves the forward
direction twice and does not prove the backward direction.
\end{pfevalenum}
\end{Answer}
\begin{Answer}{2.87}
This only proves that $4+6$ is even. It says nothing about the sum of any other two even numbers.
\end{Answer}
\begin{Answer}{2.89}
The problem is that this is actually a proof that $x+x$ is even if $x$ is even since $x=2a=y$ was assumed.
\end{Answer}
\begin{Answer}{2.90}
Notice that $4$ and $6$ are even, but $4+6=10$ is not divisible by $4$.  So clearly the statement
is incorrect.  Therefore, there must be something wrong with the proof.
The problem is the same as in the previous example--the proof assumed $x=y$, even if that was
not the intent of the writer.  So what was proven was that if $x$ is even, then $x+x$ is divisible by $4$.
\end{Answer}
\begin{Answer}{2.91}
Since it should be clear that the result ($-1=1$) is false, the proof can't possibly be correct.
\end{Answer}
\begin{Answer}{2.92}
No!  Example~\ref{minusOneIsNotOne} should have made it clear that this approach is flawed.
\end{Answer}
\begin{Answer}{2.93}
No, you should not be convinced.
As we just mentioned, whether or not the equation is true, sometimes you can work both
sides to get the same thing.  Thus the technique of working both sides is not valid.
It doesn't guarantee anything unless you already know that the equation is valid.
\end{Answer}
\begin{Answer}{2.94}
Since $p$ and $q$ are odd, we know that $p+q$ is even, and so
$\dfrac{p+q}{2}$ is an integer. But $p< q$ gives $2p < p+q < 2q $
and so $p < \dfrac{p+q}{2} < q$, that is, the average of $p$ and $q$
lies between them. Since $p$ and $q$ are consecutive primes, any
number between them is composite, and so divisible by at least two
primes. So $p+q = 2\left(\dfrac{p+q}{2}\right)$ is divisible by the
prime $2$ and by at least two other primes dividing
$\dfrac{p+q}{2}$.
\end{Answer}
\begin{Answer}{2.95}
\begin{pfevalenum}
\item
This is not correct.  It needs to be shown that $x^y$ can be written as $c/d$, where $c$ and $d$ are
integers with $d\neq 0$.  Ask yourself this:  Are $a^y$ and $b^y$ necessarily integers?
\item
This is not correct.  If $y=3/2$, what does it mean to multiple $x$ by itself one and a half times?
\end{pfevalenum}
\end{Answer}
\begin{Answer}{2.96}
The statement is false.  There are many counterexamples, but here is an easy one:
Let $x=2$ and $y=1/2$.  Then $x^y=2^{1/2}=\sqrt{2}$, which is irrational.
\end{Answer}
\begin{Answer}{2.97}
\begin{pfevalenum}
\item
This solution has two serious flaws.  First, we absolutely cannot assume $x$ is an integer.
The only thing we can assume about $x$ is that it is rational, and not every rational number
is an integer.
The other problem is that the writer proved the {\em inverse}, not the {\em contrapositive}.
What they needed to prove was that if $1/x$ is rational, then $x$ is rational.
So in actuality, we know is that $1/x$ is rational, not $x$.  We need to prove that $x$ is
rational based on the assumption that $1/x$ is rational.
\item
This is not really a proof.  It just takes the statement of the problem one step further.
Is the writer {\em sure} that $1/x$ can't be expressed as an integer over an integer?  Why?
There are just too many details omitted.
\item
The biggest flaw is that this is a proof of the {\em inverse} statement, not the {\em contrapositive}.
So even it the rest of the proof were correct, it would be proving the wrong thing
since the {\em inverse} and {\em contrapositive} are not equivalent.  But the rest is not
even entirely correct because the inverse statement is not quite true.
If $x=0$, then $p=0$ as well and the statement and proof falls apart for the same reason--you can't divide by $0$.
\item
This proof is {\em almost correct}.  It does correctly try to prove the contrapositive, and if it had done so
correctly, that would imply the original statement is true.
But there is one small problem:  If $a=0$ the proof
would fall apart because it would divide by $0$.  This possibility needs to be dealt with.
This is actually not too difficult to fix.  We just need to add the following sentence before the last sentence:
{\em ``Since $0\neq 1/x$ for any value of $x$, we know that $a\neq 0$.''.}
\item
This proof is correct.
\end{pfevalenum}
\end{Answer}
\begin{Answer}{2.98}
\begin{pfevalenum}
\item
As you will prove next, the statement is actually false.
Therefore the proof has to be incorrect.
But where did it go wrong?
It turns out they they tried to prove the wrong thing.  What needed to be proved was
``If $p$ is prime then $2^p-1$ is prime.''  They attempted to prove the {\em converse}
statement, which is not equivalent.
We can still learn something by evaluating their proof.
It turns out that the converse is actually true,
and the proof has a lot of correct elements.  Unfortunately, they are not put together
properly.   First of all, the proof seems to be a combination of a contradiction proof
and a proof by contrapositive.  They needed to pick one and stick with it.
Second, the arrows ($\rightarrow$) are confusing.  What do they mean?
I think they are supposed to be read as ``implies'', but a few more words are needed to
make the connections between these phrases.
Finally, the final statement is incorrect.  This does {\em not} prove that all
numbers of the form $2^p-1$ are prime when $p$ is prime.
\item
This proof is not even close.
This is a case of ``I wasn't sure how to prove it so I just said stuff that sounded good.''
You can't argue anything about the factors of
$2^p-1$ based on the factors of $2^p$.  Further, although $2^p-1$ being odd means
$2$ is not a factor, it doesn't tell us whether or not the number might have
{\em other} factors.
\end{pfevalenum}
\end{Answer}
\begin{Answer}{2.99}
Notice that $11$ is prime but that $2^{11}-1=23\cdot 89$ is not.
Therefore, not all numbers of the form $2^p-1$, where $p$ is prime, are prime.
\end{Answer}
